\documentclass[12pt]{article}
\usepackage{amsmath,amssymb}
\begin{document}
\section*{Research notes}

The quotient \(M_\beta=\Sigma/\langle G_\beta\rangle\) is the compact
flat \(1\times\beta\) Mobius band.  A realizing map \(f\) descends to an
embedding \(F:M_\beta\to\mathbb R^3\): condition (2) gives constancy on
orbits and injectivity on the quotient; compact-to-Hausdorff implies
embedding.  A locally finite \(G_\beta\)-invariant triangulation descends
to a finite triangulation, since a locally finite family of nonempty open
sets in a compact space is finite.

Definitions retained:
\begin{itemize}
\item In a clean triangle, the boundary edge is the ridge and the opposite
vertex is the apex.
\item A bend is a segment from the apex to a point of the ridge.
\end{itemize}

Local metric lemma: If an affine map on a clean triangle shortens the
ridge and lengthens the two side edges, then it lengthens every bend.  For
ridge endpoints \(A,B\), apex \(C\), \(P_s=(1-s)A+sB\),
\[
q(s)=|L(P_s)-L(C)|^2-|P_s-C|^2
\]
has
\[
q''=2(|L(A)-L(B)|^2-|A-B|^2)<0.
\]
Since \(q(0),q(1)>0\), concavity gives \(q(s)>0\).  Boundary subarcs are
shortened by subdividing at the finitely many boundary vertices and
adding the strict edgewise inequalities.

The bends form a circle naturally parametrized by the centerline
\(\{y=1/2\}/G_\beta\).  Each bend meets this centerline once, and every
centerline point lies in a unique bend; when the point lies on a common
non-boundary edge, that edge is the same limiting bend for the two
triangles.  Choosing a bottom-to-top orientation on the lifted bend circle
gives continuous unit directions and midpoints \(e(u),m(u)\), with
\[
e(u+1)=-e(u),\qquad m(u+1)=m(u).
\]
For \(0<r<1\),
\[
\Phi(u,r)=\big(e(u)\cdot e(u+r),
(m(u)-m(u+r))\cdot(e(u)\times e(u+r))\big).
\]
It extends over the two-point compactification of \(S^1\times(0,1)\).
The compactification coordinate
\[
(\sin\pi r\cos 2\pi(u+r/2),\sin\pi r\sin 2\pi(u+r/2),\cos\pi r)
\]
conjugates \((u,r)\mapsto(u+r,1-r)\) to the antipodal map, and
\(\Phi\) is odd.  Borsuk--Ulam gives two bends whose image lines are
perpendicular and coplanar, hence intersect.

Lower-bound geometry:
Choose the T-pair, label the bend not containing the line-intersection
point \(O\) in its image-interior by \(B\), and the other by \(T\).  Lift
\(T\) from \((0,0)\) to \((t,1)\), reflect \(x\mapsto -x\) if needed, so
\(t\ge0\).  Necessarily \(t<\beta\): if \(t>\beta\), the lift intersects
its \(G_\beta\)-translate in the interior; if \(t=\beta\), its endpoints
are equivalent.  Either contradicts the embedded nondegenerate image of a
bend.  Cutting along \(T\) gives boundary arcs
\[
H=\beta-t,\qquad D=\beta+t .
\]
Boundary shortening and bend lengthening imply
\[
H>\sqrt{1+t^2}.
\]
For \(B^*\), an endpoint \(v^*\) has distance \(>1\) from the line of
\(T^*\).  If \(A\in\{H,D\}\) is the boundary arc through the corresponding
endpoint, then
\[
A>\sqrt{5+t^2}.
\]
The elementary estimate used here: for base length \(L\) and perpendicular
height \(h\), the minimum of the sum of the two non-base sides is
\(\sqrt{L^2+4h^2}\), achieved over the midpoint of the base.

If \(A=H\), then \(2\beta=H+D>2\sqrt{5+t^2}>2\sqrt3\).  If \(A=D\), then
\[
2\beta>\sqrt{1+t^2}+\sqrt{5+t^2},\qquad
2\beta>2\sqrt{5+t^2}-2t.
\]
For
\[
a(t)=\sqrt{1+t^2}+\sqrt{5+t^2},\qquad b(t)=2\sqrt{5+t^2}-2t,
\]
\(a\) is increasing, \(b\) decreasing, and
\(a(1/\sqrt3)=b(1/\sqrt3)=2\sqrt3\).  Hence \(\beta>\sqrt3\).

Self-contained upper-bound construction:
Let \(\delta=\beta/3\).  For \(\beta>\sqrt3\), choose
\[
\max\{2/\sqrt3,\sqrt{1+\delta^2}\}<a<2\delta.
\]
The domain triangulation of \(M_\beta\) can be placed in the rectangle
\([0,\beta]\times[0,1]\), with \((0,y)\sim(\beta,1-y)\), as
\[
A_0=(0,0),\ B_0=(\delta,0),\ A_1=(2\delta,0),\ B_1=(\beta,0),
\]
\[
B_1=(0,1),\ A_2=(\delta,1),\ B_2=(2\delta,1),\ A_0=(\beta,1).
\]
Split the three width-\(\delta\) rectangles by diagonals \(A_0A_2\),
\(A_1A_2\), \(A_1A_0\).  Equivalently the faces are
\[
(A_i,B_i,A_{i+2}),\qquad (B_i,A_{i+1},A_{i+2}),\qquad i=0,1,2
\]
(indices mod \(3\)).  The non-boundary edge lengths are either \(1\)
(vertical edges \(B_iA_{i+2}\), in suitable lifts) or
\(\sqrt{1+\delta^2}\) (the diagonals \(A_iA_j\)).

Embedded six-triangle template:
For any equilateral \(E_0E_1E_2\) of side \(a\), take \(C_i(\varepsilon)\)
approaching the midpoint of side \(E_iE_{i+1}\) so that the six triangles
\[
(E_i,C_i,E_{i+2}),\qquad (C_i,E_{i+1},E_{i+2})
\]
embed as a polyhedral Mobius band.  It is enough to verify after an
affine normalization
\[
A_0=(0,0,0),\quad A_1=(1,0,0),\quad A_2=(0,1,0),
\]
\[
B_0=(1/2,0,\varepsilon),\quad
B_1=(1/2-\varepsilon,1/2,-2\varepsilon),\quad
B_2=(\varepsilon,1/2,0),\quad 0<\varepsilon<1/4.
\]
The abstract six-face complex is a compact surface.  Vertex links are
intervals:
\[
 \operatorname{lk}(A_i)=B_i-A_{i+2}-B_{i+1}-A_{i+1}-B_{i+2},\qquad
 \operatorname{lk}(B_i)=A_i-A_{i+2}-A_{i+1}
\]
(up to reversing the path).  It has one boundary cycle and Euler
characteristic \(6-12+6=0\), hence is a Mobius band.

Faces:
\[
\Delta_0=A_0B_0A_2,\quad \Delta_1=B_0A_1A_2,\quad
\Delta_2=A_1B_1A_0,\quad \Delta_3=B_1A_2A_0,
\]
\[
\Delta_4=A_2B_2A_1,\quad \Delta_5=B_2A_0A_1.
\]
Face normals:
\[
\begin{array}{lll}
n_0=(-\varepsilon,0,1/2),&
n_1=(\varepsilon,\varepsilon,1/2),&
n_2=(0,2\varepsilon,1/2),\\
n_3=(2\varepsilon,0,1/2-\varepsilon),&
n_4=(0,0,1/2-\varepsilon),&
n_5=(0,0,1/2).
\end{array}
\]
Edge-sharing pairs
\((0,1),(0,3),(1,4),(2,3),(2,5)\) are non-coplanar and meet exactly in
the common edge.  The edge-sharing pair \((4,5)\) is coplanar in \(z=0\),
but the third vertices are on opposite sides of the common edge \(B_2A_1\):
\[
 \det(A_1-B_2,A_2-B_2)=(1-2\varepsilon)/2>0,\qquad
 \det(A_1-B_2,A_0-B_2)=-1/2<0.
\]
Thus it also meets exactly in the common edge.

The other nine pairs are vertex-only and their two face normals are nonproportional.  The signs of one direction of the intersection
line in the two tangent-cone bases (the two edge directions from the
shared vertex) are:
\[
\begin{array}{c|c|c}
(i,j)&\Delta_i&\Delta_j\\ \hline
(0,2)&(-,+)&(-,+)\\
(0,4)&(-,0)&(-,+)\\
(0,5)&(0,-)&(-,+)\\
(1,2)&(-,+)&(+,-)\\
(1,3)&(+,-)&(-,+)\\
(1,5)&(0,+)&(+,-)\\
(2,4)&(0,-)&(+,-)\\
(3,4)&(0,-)&(-,+)\\
(3,5)&(0,+)&(+,-)
\end{array}
\]
No row has both columns nonnegative or both nonpositive (coordinatewise), so neither
direction of the intersection line lies in both tangent cones.  Thus the
model is embedded.

Length inequalities in the upper construction:
As \(\varepsilon\to0\),
\[
|E_iC_i|, |C_iE_{i+1}|\to a/2<\delta,\quad
|C_iE_{i+2}|\to(\sqrt3/2)a>1,\quad |E_iE_j|=a>\sqrt{1+\delta^2}.
\]
For small positive \(\varepsilon\) these strict inequalities persist.
Mapping \(A_i\mapsto E_i\), \(B_i\mapsto C_i\), affinely on faces,
therefore shortens every ridge and lengthens every non-boundary side.
The embedded quotient map lifts to the required \(G_\beta\)-invariant map
on \(\Sigma\).

\end{document}
